% Slide 6 — Hipótese
\begin{frame}{Hipótese central}
  \begin{center}
    \large O desempenho das políticas de reposição depende do
    \highlight{regime / perfil operacional da demanda}, e essa relação pode
    ser \highlight{aprendida} por um meta-modelo supervisionado.
  \end{center}
  % \vspace{0.55em}
  % \small\textcolor{slategray}{Regime, aqui, parte da taxonomia de Syntetos-Boylan: ADI mede intermitência e CV$^2$ mede variabilidade, formando quatro regimes macro de demanda.}
  % \vspace{0.55em}
  % \begin{enumerate}
  %   \item \textbf{Heterogeneidade genuína:} não existe política universalmente melhor.
  %   \item \textbf{Aprendibilidade:} a relação perfil~$\rightarrow$~política pode ser
  %     aprendida a partir de características operacionais.
  %   \item \textbf{Generalização:} um modelo treinado dessa forma pode
  %     recomendar políticas para novas séries.
  % \end{enumerate}
  % \note{
  %   A hipótese central se desdobra em três premissas. A primeira é que não
  %   existe uma política universalmente melhor -- diferentes séries favorecem
  %   políticas diferentes. Quando digo regime, estou usando como ponto de
  %   partida a taxonomia de Syntetos-Boylan, que combina ADI, a intermitência,
  %   e CV ao quadrado, a variabilidade, para formar quatro regimes macro. Os
  %   PODs, apresentados depois, refinam esses regimes com características
  %   operacionais adicionais. A segunda premissa é que essa relação entre
  %   características da série e desempenho relativo das políticas pode ser
  %   aprendida por um modelo supervisionado. A terceira, consequência das duas
  %   anteriores, é que esse modelo treinado pode generalizar para séries
  %   novas, reduzindo a necessidade de testar manualmente cada produto. Os
  %   resultados que vou apresentar sustentam parcialmente essas três premissas
  %   -- a validação completa da generalização ainda é trabalho futuro.
  % }
\end{frame}
